\section{Présentation du cas expérimental et architecture du système}

\subsection{Objectif et scénario de l'expérimentation}

L'objectif principal de cette expérimentation est de comparer deux approches de développement de systèmes de détection d'objets dangereux par vision par ordinateur :

\begin{enumerate}
    \item \textbf{Approche Baseline (vulnérable)} : Un modèle MobileNetV2 entraîné de manière standard via transfer learning, représentant les pratiques courantes de développement IA sans considération particulière pour la robustesse adversariale.

    \item \textbf{Approche Secured (robuste)} : Un modèle MobileNetV2 identique architecturalement, mais entraîné avec des défenses intégrant adversarial training et des mécanismes de détection d'anomalies.
\end{enumerate}

Le scénario choisi est la \textbf{détection binaire d'armes à feu}, un cas d'usage critique dans les systèmes de sécurité où une attaque adversariale pourrait permettre de faire passer un objet dangereux pour un objet sûr. Le problème est formulé comme une classification binaire : \textbf{classe ``safe''} (objets sûrs) versus \textbf{classe ``dangerous''} (armes à feu). Cette simplification permet de concentrer l'analyse sur la robustesse adversariale plutôt que sur la complexité taxonomique, tout en facilitant l'interprétation des résultats de sécurité.

% TODO: INSÉRER FIGURE - Architecture du cas expérimental (Baseline vs Secured)
% Fichier: figures/architecture_poc.png
\begin{figure}[H]
    \centering
    % \includegraphics[width=\textwidth]{figures/architecture_poc.png}
    \caption{Architecture du cas expérimental : comparaison entre le modèle Baseline et le modèle Secured}
    \label{fig:architecture_poc}
\end{figure}

Cette comparaison permettra de quantifier l'amélioration de robustesse apportée par les mécanismes de défense, d'évaluer le trade-off entre performance sur données propres et robustesse adversariale, et de valider la faisabilité pratique de déploiement de modèles robustes en production.


\subsection{Architecture globale et composants du système}

Le système implémente l'architecture de sécurisation en cinq zones proposée au Chapitre 2, structurée autour d'un pipeline complet allant de la préparation des données jusqu'au déploiement en production. La figure~\ref{fig:architecture_globale} illustre cette architecture modulaire.

% TODO: INSÉRER FIGURE - Architecture globale du système
% Fichier: figures/architecture_globale.png
% Voir: redaction/diagramme_architecture_globale.html
\begin{figure}[H]
    \centering
    % \includegraphics[width=\textwidth]{figures/architecture_globale.png}
    \caption{Architecture globale du système de détection sécurisé}
    \label{fig:architecture_globale}
\end{figure}

L'architecture se décompose en quatre couches principales qui interagissent de manière séquentielle. La \textbf{couche de données} gère le chargement et le prétraitement du dataset personnalisé, composé d'images d'armes à feu provenant de Kaggle et d'images sûres extraites de COCO. La \textbf{couche d'entraînement} implémente le pipeline d'adversarial training avec MobileNetV2, incluant la génération d'exemples adverses (FGSM, PGD) durant l'entraînement et le suivi des métriques de robustesse. La \textbf{couche de services} expose une API REST permettant l'inférence sur de nouvelles images et l'évaluation de la robustesse du modèle déployé. Enfin, la \textbf{couche de présentation} fournit une interface web pour l'interaction utilisateur et la visualisation des résultats.

Les composants clés du système, détaillés dans la figure~\ref{fig:diagramme_composants}, incluent le \textit{DataLoader} pour la préparation des données, le module \textit{AdversarialGenerator} pour la synthèse d'exemples adverses, le \textit{TrainingManager} coordonnant l'entraînement standard et robuste, le \textit{RobustnessEvaluator} testant le modèle contre différentes attaques, et le \textit{ModelServing} gérant le déploiement dans un conteneur Docker. Cette architecture modulaire garantit la séparation des responsabilités, facilite les tests unitaires, et permet la reproductibilité grâce au versionnement des hyperparamètres et configurations.

% TODO: INSÉRER FIGURE - Diagramme de composants détaillé
% Fichier: figures/diagramme_composants.png
% Voir: redaction/diagramme_composants.html
\begin{figure}[H]
    \centering
    % \includegraphics[width=0.95\textwidth]{figures/diagramme_composants.png}
    \caption{Diagramme de composants du système avec interactions entre modules}
    \label{fig:diagramme_composants}
\end{figure}


\subsection{Présentation du dataset personnalisé}

Le dataset utilisé pour cette expérimentation combine deux sources complémentaires pour créer un ensemble d'apprentissage équilibré de 1\,896 images réparties en deux classes. La classe \textbf{``dangerous''} (948 images) provient d'un dataset Kaggle spécialisé en détection d'armes à feu, incluant des pistolets et fusils photographiés sous différents angles, distances et conditions d'éclairage. La classe \textbf{``safe''} (948 images) est constituée d'images extraites du dataset COCO (Common Objects in Context), comprenant des personnes, des objets quotidiens et des scènes variées. Ce choix permet d'assurer une diversité suffisante pour éviter les faux positifs tout en maintenant un équilibre parfait (50\%/50\%) entre les classes.

Le tableau~\ref{tab:dataset_stats} détaille la répartition des données selon le protocole standard 70\%/15\%/15\% pour l'entraînement, la validation et le test. L'ensemble d'entraînement contient 1\,398 images (699 par classe), l'ensemble de validation 294 images (147 par classe), et l'ensemble de test 204 images (102 par classe). Cette distribution garantit une évaluation fiable tout en conservant suffisamment de données pour l'entraînement adversarial, qui nécessite la génération d'exemples perturbés à chaque batch.

\begin{table}[H]
    \centering
    \caption{Statistiques du dataset personnalisé}
    \label{tab:dataset_stats}
    \begin{tabular}{lcccc}
        \toprule
        \textbf{Ensemble} & \textbf{Classe ``safe''} & \textbf{Classe ``dangerous''} & \textbf{Total} & \textbf{Ratio} \\
        \midrule
        Entraînement & 699 & 699 & 1\,398 & 73.7\% \\
        Validation   & 147 & 147 & 294   & 15.5\% \\
        Test         & 102 & 102 & 204   & 10.8\% \\
        \midrule
        \textbf{Total} & \textbf{948} & \textbf{948} & \textbf{1\,896} & \textbf{100\%} \\
        \bottomrule
    \end{tabular}
\end{table}

Toutes les images sont prétraitées de manière uniforme : redimensionnement à $224 \times 224$ pixels (format d'entrée de MobileNetV2), normalisation des valeurs de pixels dans l'intervalle $[0, 1]$, et application d'augmentation de données durant l'entraînement uniquement. Les transformations d'augmentation incluent des rotations aléatoires ($\pm 20°$), des flips horizontaux, des zooms aléatoires (facteur 0.8 à 1.2), et des variations de luminosité ($\pm 20\%$). Ces augmentations permettent d'accroître artificiellement la diversité du dataset d'entraînement et d'améliorer la capacité de généralisation du modèle.

% TODO: GÉNÉRER CAPTURES D'ÉCRAN
% 1. Grille 4×4 d'exemples de chaque classe
%    Script: python scripts/visualize_dataset.py --class dangerous --grid 4x4 --output figures/dataset_dangerous.png
%    Script: python scripts/visualize_dataset.py --class safe --grid 4x4 --output figures/dataset_safe.png

\begin{figure}[H]
    \centering
    \begin{subfigure}[b]{0.48\textwidth}
        \centering
        % \includegraphics[width=\textwidth]{figures/dataset_dangerous.png}
        \caption{Exemples de la classe ``dangerous''}
        \label{fig:dataset_dangerous}
    \end{subfigure}
    \hfill
    \begin{subfigure}[b]{0.48\textwidth}
        \centering
        % \includegraphics[width=\textwidth]{figures/dataset_safe.png}
        \caption{Exemples de la classe ``safe''}
        \label{fig:dataset_safe}
    \end{subfigure}
    \caption{Exemples représentatifs du dataset personnalisé}
    \label{fig:dataset_examples}
\end{figure}

La figure~\ref{fig:dataset_examples} présente des échantillons représentatifs des deux classes. La combinaison de sources Kaggle et COCO offre plusieurs avantages : une diversité réaliste de contextes photographiques, la reconnaissance communautaire garantissant la qualité des données, un équilibrage parfait évitant tout biais d'apprentissage, et la présence d'images authentiques (non synthétiques) assurant la pertinence du modèle en conditions réelles de déploiement.